% A LaTeX template for EXECUTIVE SUMMARY of the MSc Thesis submissions to
% Politecnico di Milano (PoliMi) - School of Industrial and Information Engineering
%
% P. F. Antonietti, S. Bonetti, A. Gruttadauria, G. Mescolini, A. Zingaro
% e-mail: template-tesi-ingind@polimi.it
%
% Last Revision: October 2021
%
% Copyright 2021 Politecnico di Milano, Italy. Inc. All rights reserved.

\documentclass[11pt,a4paper]{article}

%------------------------------------------------------------------------------
%	REQUIRED PACKAGES AND  CONFIGURATIONS
%------------------------------------------------------------------------------
% PACKAGES FOR TITLES
\usepackage{titlesec}
\usepackage{color}
% PACKAGES FOR LANGUAGE AND FONT
\usepackage[utf8]{inputenc}
\usepackage[english]{babel}
\usepackage[T1]{fontenc} % Font encoding
% PACKAGES FOR IMAGES
\usepackage{graphicx}
\graphicspath{{Images/}} % Path for images' folder
\usepackage{eso-pic} % For the background picture on the title page
\usepackage{subfig} % Numbered and caption subfigures using \subfloat
\usepackage{caption} % Coloured captions
\usepackage{subcaption}

\usepackage{transparent}

% STANDARD MATH PACKAGES
\usepackage{amsmath}
\usepackage{amsthm}
\usepackage{bm}
\usepackage[overload]{empheq}  % For braced-style systems of equations

% PACKAGES FOR TABLES
\usepackage{tabularx}
\usepackage{longtable} % tables that can span several pages
\usepackage{colortbl}

% PACKAGES FOR ALGORITHMS (PSEUDO-CODE)
\usepackage{algorithm}
\usepackage{algorithmic}

% PACKAGES FOR REFERENCES & BIBLIOGRAPHY
\usepackage[colorlinks=true,linkcolor=black,anchorcolor=black,citecolor=black,filecolor=black,menucolor=black,runcolor=black,urlcolor=black]{hyperref} % Adds clickable links at references
\usepackage{cleveref}
\usepackage[square, numbers, sort&compress]{natbib} % Square brackets, citing references with numbers, citations sorted by appearance in the text and compressed
\bibliographystyle{plain} % You may use a different style adapted to your field

% PACKAGES FOR THE APPENDIX
\usepackage{appendix}

% PACKAGES FOR ITEMIZE & ENUMERATES
\usepackage{enumitem}

% OTHER PACKAGES
\usepackage{amsthm,thmtools,xcolor} % Coloured "Theorem"
\usepackage{comment} % Comment part of code
\usepackage{fancyhdr} % Fancy headers and footers
\usepackage{lipsum} % Insert dummy text
\usepackage{tcolorbox} % Create coloured boxes (e.g. the one for the key-words)
\usepackage{stfloats} % Correct position of the tables
\usepackage{multirow}
\usepackage{multicol}





%-------------------------------------------------------------------------
%	NEW COMMANDS DEFINED
%-------------------------------------------------------------------------
% EXAMPLES OF NEW COMMANDS -> here you see how to define new commands
\newcommand{\bea}{\begin{eqnarray}} % Shortcut for equation arrays
\newcommand{\eea}{\end{eqnarray}}
\newcommand{\e}[1]{\times 10^{#1}}  % Powers of 10 notation
\newcommand{\mathbbm}[1]{\text{\usefont{U}{bbm}{m}{n}#1}} % From mathbbm.sty
\newcommand{\pdev}[2]{\frac{\partial#1}{\partial#2}}
% NB: you can also override some existing commands with the keyword \renewcommand

%----------------------------------------------------------------------------
%	ADD YOUR PACKAGES (be careful of package interaction)
%----------------------------------------------------------------------------
\usepackage{amsfonts} 
\usepackage[font=footnotesize,labelfont=bf]{caption}

%----------------------------------------------------------------------------
%	ADD YOUR DEFINITIONS AND COMMANDS (be careful of existing commands)
%----------------------------------------------------------------------------
\usepackage{import}
\usepackage{xifthen}
\usepackage{pdfpages}
\usepackage{transparent}
\usepackage{wrapfig}
\usepackage{dsfont}
\usepackage{shadethm}

\parindent=0pt

\newcommand{\incfig}[1]{%
    \def\svgwidth{\columnwidth}
    \import{./Images/}{#1.pdf_tex}
}
\newcommand*{\oldepsilon}{\epsilon}
\renewcommand*{\epsilon}{\varepsilon}
\DeclareMathOperator*{\esssup}{ess\,sup}

\numberwithin{equation}{section}
%----------------------------------------------------------------------------
%	CONFIGURATION OF THEOREM ENVIRONMENTS


% \newtheoremstyle{break}
%     {\partopsep}{\topsep}%  
%     {\normalfont}{}
%     {\bfseries}{}%
%     {\newline}{}%
% \theoremstyle{break}
% \newtheorem{theorem}{Theorem}[section]
% \newtheorem{corollary}{Corollary}[section]
% \newtheorem{proposition}{Proposition}[section]
% \newtheorem{remark}{Remark}[section]
% \newtheorem{lemma}{Lemma}[section]
% \newtheorem{notation}{Notation}[section]
% \newtheorem{definition}{Definition}[section]

% \newtheorem{definition}{Definition}[section]

% \newtheorem*{remark}{Remark}

% \newtheorem{lemma}{Lemma}[section]
% %----------------------------------------------------------------------------

% Do not change Configuration_files/config.tex file unless you really know what you are doing.
% This file ends the configuration procedures (e.g. customizing commands, definition of new commands)
% Set the geometric layout of the document
\usepackage{geometry}
\geometry{
  top=3cm,
  left = 2.0cm,
  right = 2.0cm,
  bottom=2cm,
  headheight= 2cm,
  headsep= 0cm,
}
\raggedbottom

% Create color bluePoli (-> manuale grafica coordinata:  https://www.polimi.it/fileadmin/user_upload/il_Politecnico/grafica-coordinata/2015_05_11_46xy_manuale_grafica_coordinata.pdf)
\definecolor{bluePoli}{cmyk}{0.4,0.1,0,0.4}

% Custom theorem environments
\declaretheoremstyle[
  shaded={rulecolor=bluePoli!20, rulewidth=1pt, bgcolor=bluePoli!5},
  headfont=\color{bluePoli}\normalfont\bfseries,
  bodyfont=\color{black}\normalfont,
]{colored}

\captionsetup[figure]{labelfont={color=bluePoli}} % Set colour of the captions
\captionsetup[table]{labelfont={color=bluePoli}} % Set colour of the captions
\captionsetup[algorithm]{labelfont={color=bluePoli}} % Set colour of the captions

\theoremstyle{colored}
\newtheorem{theorem}{Theorem}[section]
\newtheorem{proposition}{Proposition}[section]
\newtheorem{definition}{Definition}[section]
\newtheorem*{remark}{Remark}
\newtheorem{lemma}{Lemma}[section]

% Enhances the features of the standard "table" and "tabular" environments.
\newcommand\T{\rule{0pt}{2.6ex}}
\newcommand\B{\rule[-1.2ex]{0pt}{0pt}}

% Algorithm description
\newcounter{algsubstate}
\renewcommand{\thealgsubstate}{\alph{algsubstate}}
\newenvironment{algsubstates}{
    \setcounter{algsubstate}{0}%
    \renewcommand{\STATE}{%
    \stepcounter{algsubstate}%
    \Statex {\small\thealgsubstate:}\space}
    }{}

% Custom theorem environment
\newcolumntype{L}[1]{>{\raggedright\let\newline\\\arraybackslash\hspace{0pt}}m{#1}}
\newcolumntype{C}[1]{>{\centering\let\newline\\\arraybackslash\hspace{0pt}}m{#1}}
\newcolumntype{R}[1]{>{\raggedleft\let\newline\\\arraybackslash\hspace{0pt}}m{#1}}

% Custom itemize environment
\setlist[itemize,1]{label=$\bullet$}
\setlist[itemize,2]{label=$\circ$}
\setlist[itemize,3]{label=$-$}
\setlist{nosep}

% Set separation of columns
\setlength{\columnsep}{30pt}

% Create command for background pic
\newcommand\BackgroundPic{% Adding background picture
	\put(230,358){
		\parbox[b][\paperheight]{\paperwidth}{%
			\vfill
			\centering
			\transparent{0.2}
			\includegraphics[width=0.8\paperwidth]{raggiera_polimi.eps}%
			\vfill
}}}

% Set indentation
%\setlength\parindent{0pt}

% Custom title commands
\titleformat{\section}
{\color{bluePoli}\normalfont\Large\bfseries}
{\color{bluePoli}\thesection.}{1em}{}
\titlespacing*{\section}
{0pt}{2ex}{1ex}

\titleformat{\subsection}
{\color{bluePoli}\normalfont\large\bfseries}
{\color{bluePoli}\thesubsection.}{1em}{}
\titlespacing*{\subsection}
{0pt}{2ex}{1ex}

\titleformat{\subsubsection}
{\color{bluePoli}\normalfont\normalsize\bfseries}
{\color{bluePoli}\thesubsubsection.}{1em}{}
\titlespacing*{\subsubsection}
{0pt}{2ex}{1ex}

% Custom headers and footers
\pagestyle{fancy}
\fancyhf{}

\fancyfoot{}
\fancyfoot[C]{\thepage} % page
\renewcommand{\headrulewidth}{0mm} % headrule width
\renewcommand{\footrulewidth}{0mm} % footrule width

\makeatletter
\patchcmd{\headrule}{\hrule}{\color{black}\hrule}{}{} % headrule
\patchcmd{\footrule}{\hrule}{\color{black}\hrule}{}{} % footrule
\makeatother

% -> Create the header
\chead[C]{
\centering
\begin{tcolorbox}[arc=0pt, boxrule=0pt, colback=bluePoli!60, width=\textwidth, colupper=white]
    \includegraphics[width=0.2\textwidth]{mimesis.png}
\end{tcolorbox}
}


% Insert here the info that will be displayed into your Title page
% -> title of your work
\renewcommand{\title}{\noindent A Mathematical Model for Alzheimer Disease:

from Microscopic to a Macroscopic Model using
the Two-Scale Homogenization Theory}

% -> author name and surname
\renewcommand{\author}{Andrea Bonifacio}
% -> MSc course
\newcommand\norm[1]{\lVert#1\rVert}
\newcommand{\course}{Methods and Models for Statistical Mechanics}
% -> advisor name and surname
\newcommand{\advisor}{Silvia Lorenzani}
% IF AND ONLY IF you need to modify the co-supervisors you also have to modify the file Configuration_files/title_page.tex (ONLY where it is marked)
%\newcommand{\firstcoadvisor}{Name Surname} % insert if any otherwise comment
%\newcommand{\secondcoadvisor}{Name Surname} % insert if any otherwise comment
% -> academic year
\newcommand{\YEAR}{2023-2024}

%-------------------------------------------------------------------------
%	BEGIN OF YOUR DOCUMENT
%-------------------------------------------------------------------------
\begin{document}
%-----------------------------------------------------------------------------
% TITLE PAGE
%-----------------------------------------------------------------------------
% Do not change Configuration_files/TitlePage.tex (Modify it IF AND ONLY IF you need to add or delete the Co-advisors)
% This file creates the Title Page of the document
% DO NOT REMOVE SPACES BETWEEN LINES!

%\twocolumn[{\begin{@twocolumnfalse}

\AddToShipoutPicture*{\BackgroundPic}

\hspace{-0.6cm}\includegraphics[width=0.6\textwidth]{logo_polimi_ing_indinf.eps}

\vspace{-1mm}
\fontsize{0.3cm}{0.5cm}\selectfont \bfseries \textsc{\color{bluePoli} Report}\\

\vspace{-0.2cm}
\Large{\textbf{\color{bluePoli}{\title}}}\\

\vspace{-0.2cm}
\fontsize{0.3cm}{0.5cm}\selectfont \bfseries \textsc{\color{bluePoli} \course}\\

\vspace{-0.2cm}
\fontsize{0.3cm}{0.5cm} \selectfont \bfseries Authors: \textsc{\textbf{\author}}\\

%\vspace{-0.4cm}
%\fontsize{0.3cm}{0.5cm}\selectfont \bfseries Advisor: \textsc{\textbf{\advisor}}\\

% if only ONE co-advisor is present:
%\vspace{-0.4cm}
%\fontsize{0.3cm}{0.5cm}\selectfont \bfseries Co-advisor: \textsc{\textbf{\firstcoadvisor}}\\
% if more than one co-advisors are present:
%\vspace{-0.4cm}
%\fontsize{0.3cm}{0.5cm}\selectfont \bfseries Co-advisors: \textsc{\textbf{\firstcoadvisor}}\textsc{\textbf{\secondcoadvisor}}\\

\vspace{-0.4cm}
\fontsize{0.3cm}{0.5cm}\selectfont \bfseries Academic year: \textsc{\textbf{\YEAR}}

\small \normalfont

\vspace{11pt}

\centerline{\rule{1.0\textwidth}{0.4pt}}

%\vspace{15pt}
%\end{@twocolumnfalse}}]

\thispagestyle{plain} % In order to not show the header in the first page


%%%%%%%%%%%%%%%%%%%%%%%%%%%%%%
%%     THESIS MAIN TEXT     %%
%%%%%%%%%%%%%%%%%%%%%%%%%%%%%%
In this work will be presented a mathematical model for Alzheimer's disease (AD) at the microscopic scale. The model will be derived from the Smoluchowski equation, which describes the evolution of the density of diffusing particles that coagulate in pairs. Then will be proved that the model two-scale converges to a macroscopic model asymptotically consistent with the original one. Indeed, the information given on the microscale by the non-homogeneous Neumann boundary condition is transferred into a source term appearing in the limiting (homogenized) equations. Most of this report is based off the work of \cite{Franchi_Lorenzani_2016}, with some additions from \cite{Bertsch}.

%-----------------------------------------------------------------------------
% INTRODUCTION
%-----------------------------------------------------------------------------

\section{Introduction}

Numerical simulations play a critical role in a wide array of scientific and engineering applications, providing insights into the behavior of physical systems under various conditions. Among the most prominent techniques for performing such simulations is Finite Element Modeling (FEM). FEM discretizes a continuous domain into a mesh of finite elements, allowing for the approximation of solutions to complex partial differential equations (PDEs). However, one significant drawback of FEM is its computational intensity, especially when high resolution is required for accurate results. This research aims to explore the potential of Deep Learning (DL) techniques to accelerate FEM simulations, focusing specifically on the deformation of objects subjected to external forces.

The deformation of an object under an applied force is directly tied to the object's discretization. In FEM, the object is represented by a mesh, where the resolution of the mesh—i.e., the size and number of elements—clearly impacts the accuracy and computational cost of the simulation. High-resolution meshes can capture fine details of deformation, leading to more accurate simulations, but they are computationally expensive and time-consuming. 

The main goal of this work is to study the efficacy of a method that combines both Finite Element Modeling (FEM) and DL to obtain a realistic simulation of an object in a fraction of the time that would be required by a traditional FEM simulation. The idea is to, somehow, train a DL model to have inside the information given by the refined discretization and pass them on a coarser discretization.

The idea of using DL techniques to solve scientific problem is not new. Thanks to the rise of new frameworks and libraries, such as TensorFlow and PyTorch, it is now possible to train very complex models on large datasets in a reasonable amount of time. For the problem at hand, a lot of different approaches can be found in the existing literature: a lot of them are based on the idea that the deep learning model should predict the whole dynamic of the system, for example MeshGraphNet \cite{pfaffLearningMeshBasedSimulation2021a} or its multiscale version \cite{fortunatoMultiScaleMeshGraphNets2022}, but these are just two examples of the many possible approaches \cite{jiangMeshfreeFlowNetPhysicsConstrainedDeep2020}, \cite{djeumouNeuralNetworksPhysicsInformed2022}, \cite{hanPredictingPhysicsMeshreduced2022a}. Other methods rely on solving a time independent problem, using various architectures, such as PINNs \cite{djeumouNeuralNetworksPhysicsInformed2022} or GNNs \cite{gaoPhysicsinformedGraphNeural2022}. The proposed method falls into the second category, as it will be explained in the following sections.

One interesting solution is given by \cite{Wang_Du_Coros_Thomaszewski_2024}, where the authors extend the concept of linear modes and modal dynamics \cite{Pentland_Williams_1989} to be able to handle larger deformations. The key idea is to obtain a linear approximation of the deformation, then train a network to minimize the energy of the system, so that for every modal coordinate, the network will learn a non-linear correction, effectively learning a series of non-linear modes. This allows for faster simulations by exploiting subspace dynamics, where the simulation is performed in a reduced space to ease computational costs. The authors show that 



\section{Conclusions and final remarks}
The model presented here can be related to the one introduced in \cite{Bertsch}, where the process of diffusion and agglomeration of $\mathrm{A} \beta$ is described on the macroscale by a Smoluchowski system with a source term, coupled with a kinetic-type transport equation for the distribution function of the degree of malfunctioning of neurons that keeps into account the spreading of the disease through a neuron-to-neuron prion-like transmission. At the beginning of the analysis we said that the $\psi(t,x,y)$ is a given for our model. In particular, the mathematical analysis carried out in this paper can be regarded as a formal derivation (that is, neglecting regularity issues), starting from the microscale, of the source term $\mathcal{F}$, introduced in the second equation of (2.16) in \cite{Bertsch}, in order to describe the production of $\mathrm{A} \beta$ in monomeric form by neurons. A comparison between the model Eq. \eqref{eq 16} and the second equation of (2.16) in \cite{Bertsch} allows to make the following choice of the function $\psi(t, x, y)$ appearing in \eqref{eq 16}:
\begin{equation}
  \psi(t, x, y)=C_{\mathcal{F}} \int_{0}^{1}\left(\mu_{0}+a\right)(1-a) f(x, a, t) \, {d} a g(y)
\label{eq 110}\end{equation}
with
$$
\int_{\Gamma} g(y) \, {d} \sigma(y)=\text { const.}
$$
In Eq. \eqref{eq 110}, the parameter $a \in[0,1]$ describes the degree of malfunctioning of a neuron: $a$ close to 0 stands for 'the neuron is healthy,' whereas $a$ close to 1 for 'the neuron is dead.' Given $x \in \Omega, t \geq 0$ and $a \in[0,1]$,
$$
f(x, a, t) \, {d} a
$$
indicates the fraction of neurons close to $x$ with degree of malfunctioning at time $t$ between $a$ and $a+d a$. Furthermore, the small constant $\mu_{0}>0$ in \eqref{eq 110} accounts for $\mathrm{A} \beta$ production by healthy neurons. In \cite{Bertsch}, an evolution equation for the distribution function $f$ has been proposed:
\begin{equation}
  \partial_{t} f+\partial_{a}(f v[f])=J[f],
\label{eq 111}\end{equation}
where $v=v(x, a, t)$ indicates the deterioration rate of the health state of the neurons. We assume that
\begin{equation}
  v[f]=\iint_{\Omega \times[0,1]} \mathcal{K}(x, a, y, b) f(y, b, t) \, {d} y \, {d} b .
\label{eq 112}\end{equation}
The integral term describes the possible prion-like propagation of AD through the neural pathway. Malfunctioning neighbors are harmful for a neuron's health state, while healthy ones are not:
\begin{align*}
&\mathcal{K}(x, a, y, b) \geq 0, & \forall x, y \in \Omega, a, b \in[0,1], \\
&\mathcal{K}(x, a, y, b)=0, & \text { if } a>b .
\end{align*}
The term $J[f]$ in \eqref{eq 111} accounts for the onset of AD, since it is written in terms of the probability that, in randomly chosen parts of the cerebral tissue, the degree of malfunctioning of neurons randomly jumps to higher values due to external agents or genetic factors. What prevents this homogenization results from being considered a fully rigorous derivation of the source term $\mathcal{F}$, appearing in the second equation of (2.16) in \cite{Bertsch}, is the fact that the solutions of Eq. \eqref{eq 111} do not satisfy, in general, all the regularity properties assumed on $\psi$. However, Eq. \eqref{eq 110} allows to establish a link, at least formally, between the limit model derived here and the one presented in \cite{Bertsch}, suggesting a possible choice for $\psi$, considered, in the present paper, as a generic given function.
It is worth noting that the plots of $f$, at different times, can be directly compared with medical fluorodeoxyglucose PET images \cite{Bertsch}. The numerical simulations reported in \cite{Bertsch} are in good qualitative agreement with clinical images of the disease distribution in the brain which vary from early to advanced stages.
It's good to notice that the homogenization theory works in every perforated domain so one can follow the previous arguments to homogenize every system in a random perforated domain.

%-----------------------------------------------------------------------------
% CONCLUSION
%-----------------------------------------------------------------------------


%---------------------------------------------------------------------------
%  BIBLIOGRAPHY
%---------------------------------------------------------------------------
%\newpage
% Remember to insert here only the essential bibliography of your work
\bibliography{bibliography.bib} % automatically inserted and ordered with this command

\end{document}
